\chapter{对比研究与应用案例分析}

\section{对比研究目标}
本章在前文基础上，进一步回答三个问题：
\begin{enumerate}
  \item 与常见高斯牛顿迭代相比，当前两种方法表现如何；
  \item 在不同初值质量下，哪种方法更稳健；
  \item 在典型应用场景下，模型应如何配置参数。
\end{enumerate}

\section{与高斯牛顿法的对比}
构建第三种对比算法：仅使用雅可比近似二阶信息，不显式计算二阶导。其更新形式为
\[
(\bm J^\T\bm W\bm J)\Delta\bm p=\bm J^\T\bm W\bm r.
\]
该式与本文法方程方法本质一致，但实现时不进行额外角度二阶项修正。实验发现：
\begin{itemize}
  \item 在小噪声、好初值场景，三者精度接近；
  \item 在中高噪声场景，牛顿法收敛步数更少；
  \item 在几何较差场景，引入阻尼的法方程法最稳健。
\end{itemize}

\section{初值质量分层实验}
将初值分为四个等级：
\begin{itemize}
  \item G1：与真值距离小于 50 m；
  \item G2：距离在 50--200 m；
  \item G3：距离在 200--800 m；
  \item G4：距离超过 800 m。
\end{itemize}
统计显示，G1-G2 区间两种算法均稳定，G3 开始出现少量失败，G4 需要结合阻尼和线搜索。

\section{案例一：单目标静态高精度定位}
在站位几何较优、传感器标定准确条件下，算法主要目标是达到高精度与高重复性。建议配置：
\begin{enumerate}
  \item 使用三测距交会初值；
  \item 优先牛顿法快速收敛；
  \item 输出 95\% 与 99\% 置信椭球作为质控结果。
\end{enumerate}

\section{案例二：机动目标近实时定位}
在近实时场景中，单帧迭代次数必须受限。建议策略：
\begin{itemize}
  \item 利用上一时刻估计作为当前初值；
  \item 采用法方程法并限制迭代步数；
  \item 每隔固定帧使用牛顿法做一次精化；
  \item 将结果输入滤波器进行时序平滑。
\end{itemize}

\section{案例三：观测异常场景}
在部分测站受遮挡或数据突变时，建议采用鲁棒流程：
\begin{enumerate}
  \item 先做残差门限筛选；
  \item 对异常测站降权；
  \item 采用 Huber 损失进行 IRLS；
  \item 在结果层面输出告警标记而非仅给出单一点估计。
\end{enumerate}

\section{工程部署清单}
为了将模型落地为可运行模块，建议配置如下清单：
\begin{table}[H]
\centering
\caption{工程部署关键项}
\begin{tabular}{ll}
\toprule
类别 & 建议 \\
\midrule
坐标体系 & 统一地理坐标与本地直角坐标转换定义 \\
时间同步 & 记录各传感器时间戳与时延补偿 \\
标定管理 & 维护站位与安装偏差参数版本 \\
异常监控 & 记录残差超阈事件与失败样本 \\
计算资源 & 预留批处理并行接口 \\
结果接口 & 输出位置、协方差、置信椭球参数 \\
\bottomrule
\end{tabular}
\end{table}

\section{评价指标体系}
建议建立多层评价指标：
\begin{itemize}
  \item \textbf{精度指标}：RMSE、MAE、95\% 误差半径；
  \item \textbf{稳定性指标}：失败率、平均迭代次数、异常触发率；
  \item \textbf{效率指标}：平均耗时、峰值耗时；
  \item \textbf{可解释性指标}：置信椭球体积、主轴方向变化。
\end{itemize}

\section{章节讨论}
通过对比与案例分析可见，算法优劣没有绝对结论，应结合任务目标选择：
\begin{enumerate}
  \item 追求收敛速度和高精度时，优先牛顿法；
  \item 追求稳健与实现简洁时，优先法方程法；
  \item 复杂场景建议采用“法方程粗解 + 牛顿精化”组合策略。
\end{enumerate}

\section{本章小结}
本章将方法比较与场景需求连接起来，给出了可执行的工程配置建议。相关结论可直接用于课程项目答辩中的“方法选择依据”和“应用落地方案”部分。
